\scanpage{132}
\deleted{Évident} Il faut prouver que si $\lambda\notin\mathbb K$ i.e. $a^2-4Q\ne0$ (cf lemme
plus bas), alors \deleted{il est pure} $\mathbb Q(\lambda)$ est purement imaginaire, i.e.
les $\lambda_\alpha$ sont tous non réels. \deleted{Cela signifie que les} \inserted{Comme $\lambda_\alpha$ est solution}
de l'équation $T^{\unsure{2}}-a_\alpha T+Q=0$, cela signifie que
$a_\alpha^2-4Q<0$, ce qui résulte du fait que $4Q-a^2$ est tot $\ge0$.

b') $\Rightarrow$ b) trivial, \deleted{\unsure{puisque}}

b) $\Longrightarrow$ c) trivial, car $\mathbb K(\lambda)$ est égal à $\mathbb K$ ou est une ext.
quadratique purement imaginaire de $\mathbb K$.

c) $\Longrightarrow$ a) car si $\lambda_\alpha=\sigma_\alpha(\lambda)$, on aura \deleted{\unsure{que} $\bar\lambda_\alpha$ est} \inserted{$\lambda_\alpha\in\mathbb L$, donc}
$\bar\lambda_\alpha$ est l'unique \deleted{\unsure{sa}} \unsure{conjuguée|conjugué} de $\lambda_\alpha$ sur $\mathbb K$, \deleted{\unsure{il}} donc
\[
\lambda_\alpha\bar\lambda_\alpha=N_{\mathbb L/\mathbb K}(\lambda_\alpha)=N_{\mathbb L/\mathbb K}(\lambda)=\lambda\bar\lambda=Q.
\]

\deleted{\uline{Corollaire} \inserted{Supposons $Q\in\mathbb Q^{*+}$.} Soit $L\subset\mathbb L$, \inserted{$L$ fini sur $\mathbb Q$, $\deleted{\unsure{L}}$ non \deleted{\unsure{réel tot}} réel (i.e. non tot. réel)} Alors $L$ est engendré par un
$\lambda$ tel que $\lambda\bar\lambda=Q$. Si $Q\in\mathbb Z$, on peut supposer même
que $\lambda$ est entier.}

\uline{Remarque} Soit $K$ une extension finie tot. réelle de $\mathbb Q$,
et $Q\in\mathbb Q^{\unsure{*+|+*}}$, $Q>1$. Je dis qu'il existe alors une extension quadr.
purement imaginaire $L$ de $K$, \deleted{telle que} et un
entier $i\ge\unsure{0|1}$, \deleted{tels que} tels que $\deleted{\unsure{K}}L$ soit engendré
\inserted{sur $\mathbb Q$} par un élément $\lambda$ tel que $\lambda\bar\lambda=Q^i$ i.e. $|\lambda|^2=Q^i$.
Si $Q\in\mathbb Z$, on peut $\bar{\mathrm m}$ prendre $\lambda$ entier algébrique.

Dém. Soit $a\in K$ tel que $\deleted{\unsure{L|K}}K=\mathbb Q(a)$. \inserted{On peut supposer $a$ entier.} Soit
$i$ assez grand pour que $\deleted{\unsure{a^2}}4Q^i-a^2$ soit tot. $>0$
(\deleted{\unsure{il est possi}} existe, car $Q\in\mathbb Q$, $Q>1$). Soit $\lambda$ solution
de l'équation $\lambda^2-a\lambda+Q^i=0$. Il satisfait à la
condition voulue; si $Q$ est entier, il est entier, à l'étant.
% archive sheet number: 131
